\documentclass[11pt,a4paper]{article}

\usepackage[T1]{fontenc}
\usepackage[utf8]{inputenc}
\usepackage[english]{babel}
\usepackage{geometry}
\usepackage{graphicx}
\usepackage{booktabs}
\usepackage{longtable}
\usepackage{array}
\usepackage{enumitem}
\usepackage{hyperref}
\usepackage{xcolor}
\usepackage{float}
\usepackage{caption}

\geometry{margin=1in}
\hypersetup{
  colorlinks=true,
  linkcolor=blue,
  urlcolor=blue,
  pdftitle={SPECTER Project Report},
  pdfauthor={Cyber Horizon Team}
}

\setlength{\parindent}{0pt}
\setlength{\parskip}{0.5em}

\title{SPECTER: Open Threat Intelligence Prototype\\Technical Report}
\author{Cyber Horizon Team}
\date{March 24, 2026}

\begin{document}
\maketitle
\tableofcontents
\newpage

\section{Executive Summary}
SPECTER is an operational Threat Intelligence prototype built for the Open Threat Intelligence Challenge. The system transforms public OSINT into structured intelligence through an end-to-end pipeline:
\begin{enumerate}[leftmargin=1.5em]
  \item Collection from multiple public providers.
  \item Normalization and de-duplication.
  \item Validation with poisoning-aware checks.
  \item Scoring and prioritization.
  \item API and dashboard delivery for SOC-style workflows.
  \item Adversarial red-blue simulation to evaluate detection robustness.
\end{enumerate}

The implementation uses a Go core for ingestion, validation, scoring, persistence, and exports, plus a Python FastAPI agent service for blue analysis, red simulation, and mirror-based adversarial metrics. A React + Tauri interface provides analyst visualization.

\section{Challenge Context and Specifications}
\subsection{Challenge Scope}
The challenge requires a functional prototype that converts OSINT into structured intelligence consumable by cybersecurity teams. Core required capabilities include:
\begin{itemize}[leftmargin=1.5em]
  \item OSINT collection from public sources.
  \item Traceability of data origins.
  \item Enrichment and scoring of intelligence signals.
  \item Noise reduction and inconsistency handling.
  \item Basic poisoning detection.
  \item SOC-oriented reporting outputs.
\end{itemize}

\subsection{Deliverables Required by the Challenge}
\begin{itemize}[leftmargin=1.5em]
  \item Presentation of problem, architecture, and use case.
  \item Video demonstration (up to 2 minutes).
  \item Written report with design, technologies, and limitations.
  \item Live or recorded prototype demo.
\end{itemize}

\subsection{Scoring Criteria (Challenge Definition)}
\begin{itemize}[leftmargin=1.5em]
  \item Technical: 12 points.
  \item Non-technical: 6 points.
  \item Bonus: 2 points.
  \item Total: 20 points.
\end{itemize}

\section{Problem Statement}
OSINT is widely available but often fragmented and difficult to operationalize directly in SOC workflows. Many demonstrations remain conceptual and do not provide deterministic processing, traceability, and analyst-ready output artifacts. SPECTER addresses this gap by implementing a runnable pipeline with explicit stages, API surfaces, artifact generation, and adversarial behavior simulation.

\section{Project Objectives}
\subsection{In-Scope Objectives}
\begin{itemize}[leftmargin=1.5em]
  \item Build a working pipeline from collection to intelligence output.
  \item Preserve source metadata and evidence traceability.
  \item Provide deterministic scoring and stage transitions.
  \item Expose machine-consumable API endpoints for operations.
  \item Generate exportable artifacts (STIX and report).
  \item Demonstrate resilience through red-blue simulation.
\end{itemize}

\subsection{Out-of-Scope for Current Iteration}
\begin{itemize}[leftmargin=1.5em]
  \item Cloud-native production deployment templates.
  \item Full enterprise IAM and RBAC model.
  \item Distributed high-scale message-queue architecture.
  \item Extensive third-party SIEM connector catalog.
\end{itemize}

\section{System Architecture}
\subsection{Architecture Overview}
SPECTER architecture consists of:
\begin{itemize}[leftmargin=1.5em]
  \item \textbf{Data sources}: crt.sh, Shodan, URLHaus, AbuseIPDB, OTX.
  \item \textbf{Go collector}: collection, normalization, de-duplication, validation, scoring.
  \item \textbf{SQLite persistence}: canonical threat record storage.
  \item \textbf{Go API}: events, metrics, exports, manual injection.
  \item \textbf{Go worker}: validated-to-scored processing loop.
  \item \textbf{Python FastAPI agent service}: blue analyst, red injector, detector context over adversarial mirror.
  \item \textbf{React + Tauri UI}: analyst-facing dashboard and export view.
\end{itemize}

\subsection{Architecture Figure}
\begin{figure}[H]
\centering
% The report tries common filenames for the architecture image.
\IfFileExists{architecture.png}{\includegraphics[width=\textwidth]{architecture.png}}{
\IfFileExists{project_architecture.png}{\includegraphics[width=\textwidth]{project_architecture.png}}{
\IfFileExists{docs/architecture.png}{\includegraphics[width=\textwidth]{docs/architecture.png}}{
\IfFileExists{docs/project_architecture.png}{\includegraphics[width=\textwidth]{docs/project_architecture.png}}{
\fbox{\parbox{0.95\textwidth}{\centering
Architecture image placeholder.\\
Place the provided architecture screenshot as one of:\\
architecture.png, project_architecture.png, docs/architecture.png, or docs/project_architecture.png
}}
}}}}
\caption{SPECTER architecture diagram provided for the project.}
\label{fig:architecture}
\end{figure}

\section{Agentic Layer Design (Blue, Red, Detector)}
\subsection{Blue Analyst Agent}
\begin{itemize}[leftmargin=1.5em]
  \item Endpoint: \texttt{/agents/blue/analyze}.
  \item Function: Summarizes threats, groups by severity, extracts top IOCs, and emits recommended actions.
  \item Value: Accelerates triage and helps SOC teams prioritize investigation.
\end{itemize}

\subsection{Red Injector Agent}
\begin{itemize}[leftmargin=1.5em]
  \item Endpoint: \texttt{/agents/red/inject}.
  \item Function: Simulates adversarial injections and poisoning behavior under controlled conditions.
  \item Value: Validates robustness and helps tune mitigation and guardrail thresholds.
\end{itemize}

\subsection{Detector Agent Context}
\begin{itemize}[leftmargin=1.5em]
  \item Implemented through poisoning-aware validation and mirror metrics context.
  \item Inputs include suspicious patterns and adversarial telemetry.
  \item Output influences quarantine signaling, detection visibility, and confidence in scored intelligence.
\end{itemize}

\subsection{Adversarial Mirror}
\begin{itemize}[leftmargin=1.5em]
  \item Supports ingestion of real IOC events and synthetic injection activity.
  \item Provides synchronized dashboard snapshot endpoints.
  \item Exposes metrics for red-blue balance and event/injection visibility.
\end{itemize}

\section{Technology Stack}
\begin{longtable}{p{0.24\textwidth} p{0.70\textwidth}}
\toprule
\textbf{Layer} & \textbf{Technology} \\
\midrule
Core services & Go 1.24, net/http, repository and pipeline modules \\
Persistence & SQLite with indexed \texttt{threat\_records} schema \\
Agentic service & Python 3.11, FastAPI, Pydantic models \\
Frontend & React 19, Vite 8, Tauri 2 \\
Build and automation & Makefile, justfile, Bash orchestration scripts \\
Testing and checks & Go test, pytest, contract tests, compile and shell syntax checks \\
\bottomrule
\end{longtable}

\section{Data Pipeline Specification}
\subsection{Processing Stages}
\begin{enumerate}[leftmargin=1.5em]
  \item \textbf{Ingested}: raw record collected from configured providers.
  \item \textbf{Validated}: deterministic detection rules applied.
  \item \textbf{Scored}: composite score computed and threat level assigned.
  \item \textbf{Quarantined}: suspicious/poisoned records isolated from normal flow.
\end{enumerate}

\subsection{Scoring Model Summary}
\begin{itemize}[leftmargin=1.5em]
  \item Baseline score starts at 20.
  \item Corroboration increases score.
  \item High-risk ports and secondary ports contribute weighted bonuses.
  \item Poison detection, suspicious rules, and synthetic indicators apply penalties.
  \item Final score is clamped to [0, 100].
  \item Threat levels map to \texttt{critical/high/medium/low}.
\end{itemize}

\section{Data Model and Persistence}
\subsection{Primary Table: \texttt{threat\_records}}
\begin{itemize}[leftmargin=1.5em]
  \item Identity and IOC: \texttt{event\_id}, \texttt{ioc\_value}, \texttt{ioc\_type}.
  \item Source traceability: \texttt{source\_name}, \texttt{source\_url}, \texttt{source\_query}, raw evidence JSON.
  \item Timing and quality: \texttt{collected\_at}, \texttt{created\_at}, \texttt{updated\_at}, corroboration count.
  \item Network features: open ports JSON, ASN.
  \item Adversarial features: synthetic flag, poison attack type, detection flags and rules.
  \item Risk outputs: composite score, threat level, days to attack, pipeline stage.
\end{itemize}

\subsection{Indexes}
\begin{itemize}[leftmargin=1.5em]
  \item Stage index for operational filtering.
  \item IOC compound index for lookup and de-duplication efficiency.
  \item Collection-time index for timeline and freshness analytics.
\end{itemize}

\section{API Specification}
\subsection{Go API Endpoints}
\begin{longtable}{p{0.11\textwidth} p{0.35\textwidth} p{0.48\textwidth}}
\toprule
\textbf{Method} & \textbf{Endpoint} & \textbf{Purpose} \\
\midrule
GET & \texttt{/health} & Liveness check for service availability. \\
GET & \texttt{/api/v1/events} & Event retrieval with stage and limit options. \\
GET & \texttt{/api/v1/events/quarantined} & Quarantine-specific event retrieval. \\
GET & \texttt{/api/v1/metrics/pipeline} & Pipeline counters and freshness telemetry. \\
POST & \texttt{/api/v1/exports/stix} & STIX artifact generation. \\
POST & \texttt{/api/v1/exports/report} & Report artifact generation. \\
POST & \texttt{/api/v1/agents/injections/trigger} & Manual adversarial injection trigger path. \\
\bottomrule
\end{longtable}

\subsection{Agent Service Endpoints}
\begin{longtable}{p{0.11\textwidth} p{0.35\textwidth} p{0.48\textwidth}}
\toprule
\textbf{Method} & \textbf{Endpoint} & \textbf{Purpose} \\
\midrule
GET & \texttt{/health} & Agent service liveness. \\
GET & \texttt{/ready} & Readiness with Go API dependency check. \\
POST & \texttt{/agents/blue/analyze} & Blue analyst workflow execution. \\
POST & \texttt{/agents/red/inject} & Red injector workflow execution. \\
POST & \texttt{/agents/run} & Generic named-agent execution. \\
POST & \texttt{/mirror/ingest} & Mirror ingest of real IOC events. \\
GET & \texttt{/mirror/events} & Mirror event feed retrieval. \\
GET & \texttt{/mirror/injections} & Injection activity retrieval. \\
GET & \texttt{/mirror/metrics} & Mirror operational metrics. \\
GET & \texttt{/mirror/dashboard} & Unified dashboard snapshot. \\
POST & \texttt{/mirror/exports/stix} & Snapshot-based STIX export. \\
POST & \texttt{/mirror/exports/report} & Snapshot-based report export. \\
\bottomrule
\end{longtable}

\section{Runtime and Deployment Specifications}
\subsection{Prerequisites}
\begin{itemize}[leftmargin=1.5em]
  \item Go 1.24+.
  \item Python 3.11+.
  \item Node.js 18+.
  \item Rust toolchain for Tauri packaging.
  \item Bash-compatible shell.
\end{itemize}

\subsection{Default Runtime Ports}
\begin{itemize}[leftmargin=1.5em]
  \item Go API: 8080.
  \item Agent service: 8001.
  \item Dashboard dev server: 5173.
\end{itemize}

\subsection{Important Configuration Variables}
\begin{itemize}[leftmargin=1.5em]
  \item Pipeline: \texttt{COLLECTION\_INTERVAL\_SECONDS}, \texttt{WORKER\_CONCURRENCY}, \texttt{DB\_DSN}.
  \item Providers: \texttt{SHODAN\_TARGETS}, \texttt{ABUSEIPDB\_TARGETS}, \texttt{OTX\_TARGETS}, \texttt{URLHAUS\_HOSTS}, \texttt{CRTSH\_QUERY}.
  \item Toggles: \texttt{ENABLE\_SHODAN}, \texttt{ENABLE\_ABUSEIPDB}, \texttt{ENABLE\_OTX}, \texttt{ENABLE\_URLHAUS}, \texttt{ENABLE\_CRTSH}.
  \item Agentic controls: \texttt{RED\_AGENT\_INTERVAL\_SECONDS}, \texttt{RED\_MAX\_RATIO}, \texttt{MIN\_REAL\_EVENTS\_BEFORE\_AUTO\_RED}.
  \item Sync controls: \texttt{GO\_SYNC\_INTERVAL\_SECONDS}, \texttt{GO\_SYNC\_BATCH\_LIMIT}, \texttt{GO\_SYNC\_ON\_STARTUP}.
\end{itemize}

\section{Reliability, Security, and Governance}
\subsection{Reliability Controls}
\begin{itemize}[leftmargin=1.5em]
  \item Provider error classification and cooldown policy.
  \item Transient backoff and rate-limit suppression behavior.
  \item Graceful service shutdown for worker process.
  \item Health/readiness endpoints for observability.
\end{itemize}

\subsection{Security and Integrity Controls}
\begin{itemize}[leftmargin=1.5em]
  \item CORS allowlist for API surfaces.
  \item Poisoning-aware validation and quarantine stage.
  \item Red/blue balance controls to prevent synthetic dominance.
  \item Controlled use of public OSINT only, aligned with challenge ethics.
\end{itemize}

\section{Testing and Validation}
\subsection{Quality Strategy}
\begin{itemize}[leftmargin=1.5em]
  \item Go tests for core modules and API behavior.
  \item Python tests for agent chains, mirror behavior, and exports.
  \item Contract tests for cross-service compatibility.
  \item Python compile checks and shell script syntax checks.
\end{itemize}

\subsection{CI-Equivalent Local Validation}
The \texttt{scripts/ci\_check.sh} profile validates:
\begin{enumerate}[leftmargin=1.5em]
  \item Go tests.
  \item Python tests.
  \item Contract-only tests.
  \item Python syntax compile checks.
  \item Shell script syntax checks.
\end{enumerate}

\section{Challenge Compliance Matrix}
\begin{longtable}{p{0.38\textwidth} p{0.56\textwidth}}
\toprule
\textbf{Challenge Requirement} & \textbf{SPECTER Implementation} \\
\midrule
Collect OSINT from public sources & Multi-provider collectors (crt.sh, Shodan, URLHaus, AbuseIPDB, OTX). \\
Document and trace data origins & Source name, query, URL, timestamps, and evidence fields stored per record. \\
Enrich and score intelligence signals & Validation rules, feature extraction, and deterministic composite scoring pipeline. \\
Filter noise and inconsistent data & De-duplication, validation checks, stage transitions, quarantine flow. \\
Basic poisoning detection & Poison-aware detection logic and red injection simulation for stress tests. \\
Produce structured TI reports & API-driven report export and STIX export artifact generation. \\
Working prototype (not concept-only) & Runnable local stack with API, worker, collector, agents, dashboard, and scripts. \\
\bottomrule
\end{longtable}

\section{Known Limitations}
\begin{itemize}[leftmargin=1.5em]
  \item SQLite persistence limits horizontal scalability in current form.
  \item External API rate limits and account constraints affect source throughput.
  \item Local-first deployment model requires further cloud-hardening work.
  \item Interoperability validation against external STIX tooling remains an ongoing task.
\end{itemize}

\section{Roadmap}
\begin{enumerate}[leftmargin=1.5em]
  \item Expand full end-to-end deterministic integration tests.
  \item Add fault-injection tests for timeout, malformed payload, and rate-limit paths.
  \item Validate STIX artifacts against external validators and import workflows.
  \item Add hosted CI workflow parity with local \texttt{ci\_check}.
  \item Finalize submission package: presentation, demo video, screenshots, and checklist.
\end{enumerate}

\section{Conclusion}
SPECTER satisfies the challenge objective by delivering a functional Threat Intelligence prototype that ingests public OSINT, performs traceable and poisoning-aware processing, and outputs operational intelligence via APIs, dashboard views, and export artifacts. The architecture is modular, demonstrable, and ready for further hardening and scale-oriented evolution.

\appendix
\section{Build and Run Quick Reference}
\begin{itemize}[leftmargin=1.5em]
  \item Bootstrap: \texttt{./scripts/bootstrap.sh}
  \item Start local stack: \texttt{./scripts/run\_local.sh start}
  \item Status: \texttt{./scripts/run\_local.sh status}
  \item Logs: \texttt{./scripts/run\_local.sh logs}
  \item Stop: \texttt{./scripts/run\_local.sh stop}
  \item CI-equivalent checks: \texttt{./scripts/ci\_check.sh}
\end{itemize}

\section{LaTeX Compilation Command}
Compile from the repository root:
\begin{verbatim}
pdflatex -output-directory=docs docs/SPECTER_Project_Report.tex
\end{verbatim}

\end{document}
